\section{Introduction}
\label{sec:introduction}

Type~2 diabetes mellitus (T2D) is a chronic metabolic disorder in which the body progressively loses its ability to regulate blood glucose, leading to serious complications including cardiovascular disease, chronic kidney disease, peripheral neuropathy, and retinopathy. The International Diabetes Federation estimates that 537~million adults worldwide are living with diabetes, and that 44.7\% of all cases remain undiagnosed~\cite{Sun2022idf, IDF2021news}. Routine health checkups already capture many of the indicators relevant to T2D risk---body mass index, blood pressure, liver enzymes, and lipid profiles---yet clinicians reviewing individual test results in isolation can miss the subtle multi-variable patterns that collectively signal elevated risk. The U.S.\ Preventive Services Task Force recommends screening adults aged 35--70 who have overweight or obesity, but questionnaire-based risk scores such as FINDRISC and IDRS have demonstrated limited sensitivity in diverse populations~\cite{USPSTF2021screening, Doddamani2021comparative}. A more systematic approach to exploiting the data already collected in standard health examinations could substantially improve early detection.

Machine learning (ML) offers a principled framework for discovering high-dimensional, nonlinear relationships in routine clinical data~\cite{Rajkomar2019ml}, and recent bibliometric analyses confirm rapidly growing interest in ML-based T2D prediction~\cite{Kiran2025bibliometric}. The Korean National Health Insurance Service (NHIS) General Health Examination dataset provides approximately one million anonymized records from 2024, each containing demographics, anthropometrics, blood pressure, liver function tests, and---for a subset of examinees---a full lipid panel~\cite{NHIS2025dataset, Lim2025nhis}. This dataset is among the largest publicly available collections of real-world health screening data and mirrors the structure of clinical examinations conducted worldwide. Prior work on NHIS data has applied deep learning to diabetes prediction but has not systematically compared gradient-boosted tree ensembles against neural network architectures under identical preprocessing and evaluation conditions~\cite{Rhee2021deeplearning}. Addressing this gap is important because recent evidence suggests that tree-based models often match or exceed deep learning on tabular data, yet most clinical studies do not perform controlled head-to-head comparisons.

The contributions of this paper are as follows:
\begin{enumerate}
    \item A three-tier feature ablation study quantifying the marginal predictive value of laboratory tests and lipid panel data beyond basic anthropometric and demographic features.
    \item An eight-model comparison spanning gradient-boosted trees, neural networks, support vector machines, and logistic regression, trained and evaluated on identical 60/20/20 stratified splits of approximately one million records.
    \item A clinical threshold tuning and Platt calibration pipeline designed to produce well-calibrated risk probabilities suitable for screening deployment.
    \item A SHAP-based interpretability analysis confirming that the most influential model features align with established clinical risk factors for T2D.
    \item A deployed web-based screening tool that maps calibrated model outputs to a four-tier risk stratification for clinical use.
\end{enumerate}
The remainder of this paper is organized as follows. Section~\ref{sec:related-work} reviews related work in ML-based diabetes prediction. Section~\ref{sec:dataset} describes the NHIS dataset and preprocessing pipeline. Section~\ref{sec:methodology} details the modeling methodology. Section~\ref{sec:results} presents experimental results. Section~\ref{sec:discussion} interprets the findings and compares them with prior work. Section~\ref{sec:limitations} acknowledges limitations, and Section~\ref{sec:conclusion} concludes with future directions.
